%% =========================================================================
% Colores personalizados (debe cargarse antes de tikz y titlesec)
\usepackage[dvipsnames,svgnames,x11names]{xcolor}
%% preamble.tex - Configuración Estética y Paquetes Avanzados
%% =========================================================================

% 1. Codificación e Idioma
\usepackage[utf8]{inputenc}
\usepackage[T1]{fontenc}
\usepackage[spanish,es-tabla]{babel} % es-tabla fuerza "Tabla" en vez de "Cuadro"

% 2. Geometría de Página Académica Elegante
\usepackage{geometry}
\geometry{
    a4paper,
    total={160mm,240mm},
    left=30mm,
    top=25mm,
    headheight=15pt
}

% 3. Matemáticas Rigurosas y Avanzadas
\usepackage{amsmath}
\usepackage{amssymb}
\usepackage{amsfonts}
\usepackage{mathtools}
\usepackage{bm} % Negritas matemáticas para vectores/tensores

% 4. Química Estructural Vectorial
\usepackage{chemfig}
\usepackage[version=4]{mhchem}
\setchemfig{atom sep=1.8em, double bond sep=2.5pt, bond offset=1.5pt}

% 5. Figuras y Gráficos Nativos
\usepackage{graphicx}
\usepackage{tikz}
\usetikzlibrary{shapes,arrows,positioning,calc}
\usepackage{pgfplots}
\pgfplotsset{compat=1.18}

% 6. Estilos de Tablas de Alto Impacto Editorial
\usepackage{booktabs}        % Elimina las líneas divisorias verticales feas
\usepackage{array}           % Alineaciones avanzadas de celdas
\usepackage{siunitx}        % Alineación numérica por punto decimal y notación científica
\sisetup{output-decimal-marker = {.}}
\usepackage{threeparttable} % Notas al pie en tablas
\usepackage{multirow}       % Combinación de filas

% 7. Cabeceras y Pies de Página Premium
\usepackage{fancyhdr}
\pagestyle{fancy}
\fancyhf{}
\fancyhead[LE,RO]{\thepage}
\fancyhead[RE]{\leftmark}
\fancyhead[LO]{\rightmark}
\renewcommand{\headrulewidth}{0.4pt}
\renewcommand{\footrulewidth}{0pt}

% 8. Hipervínculos e Integración Digital
\usepackage{hyperref}
\hypersetup{
    colorlinks=true,
    linkcolor=teal,
    citecolor=blue,
    filecolor=magenta,      
    urlcolor=cyan,
    pdftitle={Tratado de Ciencia e Ingenieria de Polimeros},
    pdfauthor={Placeholder Autor},
    pdfpagemode=FullScreen,
}

% 9. Ajustes Tipográficos de Precisión (Microtype)
% expansion=false evita el error con fuentes no escalables en pdflatex
\usepackage[protrusion=true,expansion=false]{microtype}

% Configuración del formato de títulos
\usepackage{titlesec}
\titleformat{\chapter}[display]
  {\normalfont\huge\bfseries\color{teal}}
  {\chaptercomputermodern\ {\fontsize{60}{60}\selectfont\thechapter}}
  {20pt}
  {\Huge}
\newcommand{\chaptercomputermodern}{\scshape\Large\chaptername}

\titleformat{\section}
  {\normalfont\Large\bfseries\color{teal}}{\thesection}{1em}{}
\titleformat{\subsection}
  {\normalfont\large\bfseries\color{black!80}}{\thesubsection}{1em}{}
